\section{Extended deformations of dimer models}\label{sec_def_exdeform}

In this section, we add some procedures to the deformations given in Definitions~\ref{def_deformation_zig} and \ref{def_deformation_zag}
to construct a consistent dimer model having the same properties as Theorem~\ref{mutation=deformation1} in a more general situation.
The operations which will be introduced in this section might be complicated, thus we refer the reader to Figures~\ref{def_leftright}--\ref{fig_observation_zigzag4}
in the next section and Appendix~\ref{app_large_example} for recognizing their points.

\begin{setting}\label{def_deformation_exdata}
Let $\Gamma$ be a reduced consistent dimer model.
As in Definition~\ref{def_deformation_data}, we choose a~type~I zigzag path $z$, and we let $2n\coloneqq\ell(z)$ and $v\coloneqq [z]$.
We fix positive integers $r, h$ such that $r\le \big|\calZ_v^\rmI(\Gamma)\big|$ and $n=r+h$,
and take a subset $\{z_1,\dots,z_r\}\subset\calZ_v^\rmI(\Gamma)$ of type I zigzag paths.
We recall that we have $2n=\ell(z_1)=\cdots=\ell(z_r)$ by Lemma~\ref{lem_same_length}, and hence we write $z_i$ as
\begin{displaymath}
z_i=z_i[1]z_i[2]\cdots z_i[2n-1]z_i[2n].
\end{displaymath}

Then we prepare the additional data as follows.
\begin{itemize}\itemsep=0pt
\item[(1)] We consider all zigzag paths $x_1,\dots, x_s$ (resp.~$y_1,\dots, y_t$) intersecting with $z$ at some zags (resp.\ zigs) of $z$.
In this case, each of $x_1,\dots, x_s$ (resp.\ $y_1,\dots, y_t$) intersects with any $z_i$ at some zags (resp.\ zigs) of $z_i$ for all $i=1,\dots,r$ by Lemma~\ref{intersect_zigorzag2}.
We may assume that $z_1,\dots,z_r$ are ordered cyclically in the sense that if $x_j$ (resp.~$y_k$) intersects with $z_i$,
then it also intersects with $z_{i-1}$ (resp.~$z_{i+1}$).
\item[(2)]
We recall that $|x_j\cap z_i|$ (resp.~$|y_k\cap z_i|$) is the same number for any $i=1,\dots,r$ by Lemma~\ref{intersect_zigorzag2}.
Thus, for $j\in\{1,\dots, s\}$ let $m_j\coloneqq|x_j\cap z_i|$ be this constant, and for $k\in\{1,\dots, t\}$ let $m^\prime_k\coloneqq|y_k\cap z_i|$.
We note that $n=m_1+\cdots+m_s=m^\prime_1+\cdots+m^\prime_t$.
\item[(3)] Then, we divide each zigzag path $x_j$ into $m_j$ parts $x_j^{(1)},\dots,x_j^{(m_j)}$ as follows.
We first fix one of the intersections of $z_r$ and $x_j$ as the starting edge of $x_j^{(1)}$,
and tracing along $x_j$ we will arrive at another intersection of~$z_r$ and~$x_j$.
We consider the edge of $x_j$ just before this intersection as the ending edge of $x_j^{(1)}$, and consider the second intersection as the starting edge of $x_j^{(2)}$.
Repeating this procedure, we obtain $x_j^{(1)},\dots,x_j^{(m_j)}$ and the set of sub-zigzag paths:
\begin{equation}\label{subzigzag_x}
\big\{x_1^{(1)},\dots,x_1^{(m_1)}, x_2^{(1)},\dots,x_2^{(m_2)}, \dots, x_s^{(1)},\dots,x_s^{(m_s)}\big\}.
\end{equation}

Similarly, we also divide each zigzag path $y_k$ into $m^\prime_k$ parts $y_k^{(1)},\dots,y_k^{(m^\prime_k)}$
by considering one of the intersections of $z_1$ and $y_k$ as the starting edge of $y_k^{(1)}$, and obtain the set of sub-zigzag paths:
\begin{equation}\label{subzigzag_y}
\big\{y_1^{(1)},\dots,y_1^{(m^\prime_1)}, y_2^{(1)},\dots,y_2^{(m^\prime_2)}, \dots, y_t^{(1)},\dots,y_t^{(m^\prime_t)}\big\}.
\end{equation}
\item[(4)]
We then assign one of $\{z_1, \dots, z_r\}$ to $x_j^{(a_j)}$ for $j=1, \dots, s$ and $a_j=1, \dots, m_j$.
Then, we define $X_i$ as the set of edges consisting of the intersections between $z_i$ and the sub-zigzag paths in~(\ref{subzigzag_x}) that are assigned with~$z_i$.
We can make this assignment so that $|X_i|\ge1$, and then set $p_i\coloneqq|X_i|-1$ for $i=1,\dots,r$.
We call $\calX\coloneqq\{X_1,\dots,X_r\}$ a \emph{zig-deformation parameter} with respect to $z_1,\dots,z_r$ and call the tuple $\bfp=(p_1,\dots,p_r)\in\ZZ^r_{\ge0}$ the \emph{weight} of $\calX$.
Similarly, we assign one of $\{z_1, \dots, z_r\}$ to $y_k^{(b_k)}$ for $k=1, \dots, t$ and $b_k=1, \dots, m^\prime_k$.
Then, we define $Y_i$ as the set of edges consisting of the intersections between $z_i$ and the sub-zigzag paths in~(\ref{subzigzag_y}) that are assigned with~$z_i$.
We can make this assignment so that $|Y_i|\ge1$, and then set $q_i\coloneqq|Y_i|-1$ for $i=1,\dots,r$.
We call $\calY\coloneqq\{Y_1,\dots,Y_r\}$ a \emph{zag-deformation parameter} with respect to $z_1,\dots,z_r$ and call the tuple $\bfq=(q_1,\dots,q_r)\in\ZZ^r_{\ge0}$ the \emph{weight} of $\calY$.
We remark that $p_1+\cdots+p_r=m_1+\cdots+m_s-r=n-r=h$, and also that $q_1+\cdots+q_r=h$.
\end{itemize}
\end{setting}

\begin{Definition}[extended zig-deformation]\label{def_exdeformation_zig}
Let the notation be the same as in Setting~\ref{def_deformation_exdata}.
For a zig-deformation parameter $\calX=\{X_1,\dots,X_r\}$ of weight $\bfp=(p_1,\dots,p_r)$, we consider
the operations (zig-1), (zig-2), (zig-3) defined in Definition~\ref{def_deformation_zig} and use the same notation.
Then we conduct the following procedures:
\begin{itemize}\itemsep=0pt
\setlength{\parskip}{0pt}
\setlength{\leftskip}{0.3cm}
\item[(\text{zig}-4)]
For $m=1,\dots,n$ and $i=1, \dots, r$, if the zag $z_i[2m]$ of the original zigzag path~$z_i$ on~$\Gamma$ is not contained in~$X_i$,
then we add edges, which we call \emph{bypasses}
(since these edges provide a new route that connect edges of the zigzag path $x_j$ contained in non-deformed parts, see Figure~\ref{fig_observation_zigzag4}),
connecting the following pairs of black and white nodes:
\begin{align*}
(w_{i,1}[2m-1], &\, b_i[2m+1]), \,(w_{i,2}[2m-1],b_{i,1}[2m+1]), \\ &\dots, (w_{i,p_i}[2m-1],b_{i,p_i-1}[2m+1]),\, (w_i[2m-1],b_{i,p_i}[2m+1]).
\end{align*}
We denote the resulting dimer mode by $\overline{\nu}^\zig_\calX(\Gamma, \{z_1,\dots,z_r\})$.

We note that $\overline{\nu}^\zig_\calX(\Gamma, \{z_1,\dots,z_r\})$ is non-degenerate by Proposition~\ref{prop_nondegenerate}.
\item[(\text{zig}-5)]
Then, we make the dimer model $\overline{\nu}^\zig_\calX(\Gamma, \{z_1,\dots,z_r\})$ consistent using the method given in the proof of \cite[Theorem~1.1]{BIU} (see Operation~\ref{operation_BIU} and Proposition~\ref{prop_make_consistent}).
\end{itemize}

At the end, we perform the operation (join) if the resulting dimer model contains $2$-valent nodes.
We denote the resulting dimer model by $\nu^\zig_\calX(\Gamma, \{z_1,\dots,z_r\})$,
and call it the \emph{extended zig-deformation of $\Gamma$ at $\{z_1,\dots,z_r\}$ with respect to the zig-deformation parameter $\calX$}.
If the situation is clear, we simply denote this by $\nu^\zig_\calX(\Gamma)$.
\end{Definition}

\begin{figure}[h!]\centering\vspace*{-10mm}
\scalebox{0.55}{
\begin{tikzpicture}
\newcommand{\edgewidth}{0.05cm} %line_width_of _edges
\newcommand{\nodewidth}{0.05cm} %line_width_around_nodes
\newcommand{\noderad}{0.16} %radius_of_node
%coordinate setting
\coordinate (W1) at (0,1.2); \coordinate (W2) at (0,3.6);
%\coordinate (B1) at (3,0); \coordinate (B2) at (3,2.4);
\path (W1) ++(-20:4cm) coordinate (B1); \path (W2) ++(-20:4cm) coordinate (B2);
\path (B1) ++(200:2cm) coordinate (B0); \path (W2) ++(20:2cm) coordinate (W3);

\path (W1) ++(-20:0.8cm) coordinate (B1-1); \path (W1) ++(-20:1.6cm) coordinate (W1-1);
\path (W1) ++(-20:2.4cm) coordinate (B1-2); \path (W1) ++(-20:3.2cm) coordinate (W1-2);
\path (W2) ++(-20:0.8cm) coordinate (B2-1); \path (W2) ++(-20:1.6cm) coordinate (W2-1);
\path (W2) ++(-20:2.4cm) coordinate (B2-2); \path (W2) ++(-20:3.2cm) coordinate (W2-2);

\path (B1) ++(30:1cm) coordinate (B1e); \path (B1) ++(330:1cm) coordinate (B1s);
\path (W1) ++(150:1cm) coordinate (W1w); \path (W1) ++(210:1cm) coordinate (W1s);
\path (B2) ++(30:1cm) coordinate (B2e); \path (B2) ++(330:1cm) coordinate (B2s);
\path (W2) ++(150:1cm) coordinate (W2w); \path (W2) ++(210:1cm) coordinate (W2s);

\path (W1) ++(160:0.7cm) coordinate (W1+); \path (W1) ++(200:0.7cm) coordinate (W1-);
\path (W2) ++(160:0.7cm) coordinate (W2+); \path (W2) ++(200:0.7cm) coordinate (W2-);
\path (B1) ++(20:0.7cm) coordinate (B1+); \path (B1) ++(340:0.7cm) coordinate (B1-);
\path (B2) ++(20:0.7cm) coordinate (B2+); \path (B2) ++(340:0.7cm) coordinate (B2-);

\node (original) at (0,0) {
\begin{tikzpicture}
%edge
\draw [line width=\edgewidth] (B1)--(W1); \draw [line width=\edgewidth] (B2)--(W1) ; \draw [line width=\edgewidth] (B2)--(W2) ;
\draw [line width=\edgewidth] (B1)--(B0) ; \draw [line width=\edgewidth] (W2)--(W3) ;
\draw [line width=\edgewidth] (B1)--(B1e); \draw [line width=\edgewidth] (B1)--(B1s);
\draw [line width=\edgewidth] (W1)--(W1w); \draw [line width=\edgewidth] (W1)--(W1s);
\draw [line width=\edgewidth] (B2)--(B2e); \draw [line width=\edgewidth] (B2)--(B2s);
\draw [line width=\edgewidth] (W2)--(W2w); \draw [line width=\edgewidth] (W2)--(W2s);

\draw [line width=\edgewidth,line cap=round, dash pattern=on 0pt off 2.5\pgflinewidth] (W1+)--(W1-);
\draw [line width=\edgewidth,line cap=round, dash pattern=on 0pt off 2.5\pgflinewidth] (W2+)--(W2-);
\draw [line width=\edgewidth,line cap=round, dash pattern=on 0pt off 2.5\pgflinewidth] (B1+)--(B1-);
\draw [line width=\edgewidth,line cap=round, dash pattern=on 0pt off 2.5\pgflinewidth] (B2+)--(B2-);
%blacknode
\draw [line width=\nodewidth, fill=black] (B1) circle [radius=\noderad] ; \draw [line width=\nodewidth, fill=black] (B2) circle [radius=\noderad] ;
%whitenode
\draw [line width=\nodewidth, fill=white] (W1) circle [radius=\noderad] ; \draw [line width=\nodewidth, fill=white] (W2) circle [radius=\noderad] ;

%zigzag path
\coordinate (W1z) at (0,1.2); \coordinate (W2z) at (0,3.6);
\path (W1) ++(90:0.2cm) coordinate (W1z); \path (W2) ++(90:0.2cm) coordinate (W2z);
\path (W1z) ++(-20:4cm) coordinate (B1z); \path (W2z) ++(-20:4cm) coordinate (B2z);
\path (B1z) ++(200:2cm) coordinate (B0z); \path (W2z) ++(20:2cm) coordinate (W3z);
\path (B1z) ++(200:2cm) coordinate (B0z); \path (W2z) ++(20:2.2cm) coordinate (W3z);

\path (B2) ++(-90:0.2cm) coordinate (B2zz); \path (B2zz) ++(330:1cm) coordinate (B2szz);
\path (W1) ++(-90:0.2cm) coordinate (W1zz); \path (W1zz) ++(150:1.5cm) coordinate (W1wzz);

\draw [->, rounded corners, line width=0.08cm, blue] (B2szz)--(B2zz)--(W1zz)--(W1wzz);
\draw [->, rounded corners, line width=0.08cm, red] (B0z)--(B1z)--(W1z)--(B2z)--(W2z)--(W3z);

\node at (2.5,3.5) {\color{red}\large$z_i[2m+1]$};
\node at (1.2,2.2) {\color{red}\large$z_i[2m]$};
\node at (2.5,1.1) {\color{red}\large$z_i[2m-1]$};
\node at (-1.2,2.2) {\color{blue}\large$x_j$};
\end{tikzpicture}};

\node (before) at (10,0) {
\begin{tikzpicture}
%edge
\draw [line width=\edgewidth] (B1)--(W1); \draw [line width=\edgewidth] (B2)--(W2) ;
\draw [line width=\edgewidth] (B1)--(B1e); \draw [line width=\edgewidth] (B1)--(B1s);
\draw [line width=\edgewidth] (W1)--(W1w); \draw [line width=\edgewidth] (W1)--(W1s);
\draw [line width=\edgewidth] (B2)--(B2e); \draw [line width=\edgewidth] (B2)--(B2s);
\draw [line width=\edgewidth] (W2)--(W2w); \draw [line width=\edgewidth] (W2)--(W2s);

\draw [line width=\edgewidth,line cap=round, dash pattern=on 0pt off 2.5\pgflinewidth] (W1+)--(W1-);
\draw [line width=\edgewidth,line cap=round, dash pattern=on 0pt off 2.5\pgflinewidth] (W2+)--(W2-);
\draw [line width=\edgewidth,line cap=round, dash pattern=on 0pt off 2.5\pgflinewidth] (B1+)--(B1-);
\draw [line width=\edgewidth,line cap=round, dash pattern=on 0pt off 2.5\pgflinewidth] (B2+)--(B2-);

\draw [line width=\edgewidth] (B2-1)--(W1-1); \draw [line width=\edgewidth] (B2-2)--(W1-2);
\path (B1-1) ++(-70:1.5cm) coordinate (B1-1-); \draw [line width=\edgewidth] (B1-1)--(B1-1-);
\path (B1-2) ++(-70:1.5cm) coordinate (B1-2-); \draw [line width=\edgewidth] (B1-2)--(B1-2-);
\path (W2-1) ++(110:1.5cm) coordinate (W2-1-); \draw [line width=\edgewidth] (W2-1)--(W2-1-);
\path (W2-2) ++(110:1.5cm) coordinate (W2-2-); \draw [line width=\edgewidth] (W2-2)--(W2-2-);
%blacknode
\draw [line width=\nodewidth, fill=black] (B1) circle [radius=\noderad] ; \draw [line width=\nodewidth, fill=black] (B2) circle [radius=\noderad] ;
\draw [line width=\nodewidth, fill=black] (B1-1) circle [radius=\noderad] ; \draw [line width=\nodewidth, fill=black] (B1-2) circle [radius=\noderad] ;
\draw [line width=\nodewidth, fill=black] (B2-1) circle [radius=\noderad] ; \draw [line width=\nodewidth, fill=black] (B2-2) circle [radius=\noderad] ;
%whitenode
\draw [line width=\nodewidth, fill=white] (W1) circle [radius=\noderad] ; \draw [line width=\nodewidth, fill=white] (W2) circle [radius=\noderad] ;
\draw [line width=\nodewidth, fill=white] (W1-1) circle [radius=\noderad] ; \draw [line width=\nodewidth, fill=white] (W1-2) circle [radius=\noderad] ;
\draw [line width=\nodewidth, fill=white] (W2-1) circle [radius=\noderad] ; \draw [line width=\nodewidth, fill=white] (W2-2) circle [radius=\noderad] ;
\end{tikzpicture}};

\node (after) at (20,0) {
\begin{tikzpicture}
%edge
\draw [line width=\edgewidth] (B1)--(W1); \draw [line width=\edgewidth] (B2)--(W2) ;
\draw [line width=\edgewidth] (B1)--(B1e); \draw [line width=\edgewidth] (B1)--(B1s);
\draw [line width=\edgewidth] (W1)--(W1w); \draw [line width=\edgewidth] (W1)--(W1s);
\draw [line width=\edgewidth] (B2)--(B2e); \draw [line width=\edgewidth] (B2)--(B2s);
\draw [line width=\edgewidth] (W2)--(W2w); \draw [line width=\edgewidth] (W2)--(W2s);

\draw [line width=\edgewidth,line cap=round, dash pattern=on 0pt off 2.5\pgflinewidth] (W1+)--(W1-);
\draw [line width=\edgewidth,line cap=round, dash pattern=on 0pt off 2.5\pgflinewidth] (W2+)--(W2-);
\draw [line width=\edgewidth,line cap=round, dash pattern=on 0pt off 2.5\pgflinewidth] (B1+)--(B1-);
\draw [line width=\edgewidth,line cap=round, dash pattern=on 0pt off 2.5\pgflinewidth] (B2+)--(B2-);

\draw [line width=\edgewidth] (B2-1)--(W1-1); \draw [line width=\edgewidth] (B2-2)--(W1-2);
\path (B1-1) ++(-70:1.5cm) coordinate (B1-1-); \draw [line width=\edgewidth] (B1-1)--(B1-1-);
\path (B1-2) ++(-70:1.5cm) coordinate (B1-2-); \draw [line width=\edgewidth] (B1-2)--(B1-2-);
\path (W2-1) ++(110:1.5cm) coordinate (W2-1-); \draw [line width=\edgewidth] (W2-1)--(W2-1-);
\path (W2-2) ++(110:1.5cm) coordinate (W2-2-); \draw [line width=\edgewidth] (W2-2)--(W2-2-);

\draw [line width=\edgewidth] (B2-1)--(W1); \draw [line width=\edgewidth] (B2-2)--(W1-1); \draw [line width=\edgewidth] (B2)--(W1-2);
%blacknode
\draw [line width=\nodewidth, fill=black] (B1) circle [radius=\noderad] ; \draw [line width=\nodewidth, fill=black] (B2) circle [radius=\noderad] ;
\draw [line width=\nodewidth, fill=black] (B1-1) circle [radius=\noderad] ; \draw [line width=\nodewidth, fill=black] (B1-2) circle [radius=\noderad] ;
\draw [line width=\nodewidth, fill=black] (B2-1) circle [radius=\noderad] ; \draw [line width=\nodewidth, fill=black] (B2-2) circle [radius=\noderad] ;
%whitenode
\draw [line width=\nodewidth, fill=white] (W1) circle [radius=\noderad] ; \draw [line width=\nodewidth, fill=white] (W2) circle [radius=\noderad] ;
\draw [line width=\nodewidth, fill=white] (W1-1) circle [radius=\noderad] ; \draw [line width=\nodewidth, fill=white] (W1-2) circle [radius=\noderad] ;
\draw [line width=\nodewidth, fill=white] (W2-1) circle [radius=\noderad] ; \draw [line width=\nodewidth, fill=white] (W2-2) circle [radius=\noderad] ;
\end{tikzpicture}};

\path (original) ++(0:3.5cm) coordinate (original+); \path (before) ++(180:3.5cm) coordinate (before+);
\draw [->, line width=\edgewidth] (original+)--(before+) node[midway,xshift=-0.5cm,yshift=1.2cm] {\LARGE(zig-1)}
node[midway,xshift=0.5cm,yshift=0.5cm] {\LARGE--\,(zig-3)};

\path (before) ++(0:3.5cm) coordinate (before-); \path (after) ++(180:3.5cm) coordinate (after+);
\draw [->, line width=\edgewidth] (before-)--(after+) node[midway,xshift=0cm,yshift=0.5cm] {\LARGE(zig-4)} ;

\end{tikzpicture}
}
\caption{The operations (zig-1)--(zig-4) at $z_i$ for the case $p_i=|X_i|-1=2$ and $z_i[m]\not\in X_i$.}\label{fig_exzigdeform_p2}
\end{figure}
